\documentclass{article}

\usepackage[T1]{fontenc}
\usepackage[utf8]{inputenc}
\usepackage[english, russian]{babel}

\usepackage{amsmath}

\newcommand{\rl}{\hspace*{5mm}}
\renewcommand{\phi}{\varphi}

\title{Падарунак маяму шчасьцю}    
\author{Твой Гошка\thanks{Я цябе кахаю!}}
\date{11 жніўня 2023 г.}
\begin{document}
\maketitle
\begin{center}
    Прыветаньне, маё кацянё! Я зрабіў цябе падарунак на твой дзень нараджэньня, ты зараз чытаеш тлумачэньне, як працуе то, што ты ўбачыла. \\
    \vspace{3mm}
    Я цябе вельмі кахаю і жадаю, каб ты заставалася самай прыгожай, самай добрай і найлепшай дзяўчынкай, і каб ты, маленькі цуд, працягваў радаваць увесь сьвет сабою! Ты ў мяне такое сонейка!
\end{center}
\newpage
\section{Вектара}
\rl Як ты магла ўбачыць, гэта ў сваёй сутнасьці набор вектароў, якія кружацца па сваёй асі.
Любы вектар можна прадставіць у выглядзе лічб, у нашым выпадку гэта будуць камплексныя.
Нагадаю, што камплексная лічба \--- гэта пара $(a, b)$, дзе $a$ \--- рэальная частка, а $b$ \--- мнімая, і сама гэтая лічба $a + bi$,
а $i$ вызначаецца наступным чынам: $i^2 = -1$
\[ \begin{pmatrix} x \\ y \end{pmatrix} \leftrightarrow x + yi \]
\begin{center}
    \scriptsize Сувязь паміж вектарам $\begin{pmatrix} x \\ y \end{pmatrix}$ і камплекснай лічбай $x + yi$: можна правесьці біекцыю
\end{center}
\section{Формула Эйлера}
\rl Як мы ведаем, экспаненту і трыганаметрычныя функцыі магчыма раскласьці ў рад Тэйлара
\[ e^x = \sum_{n = 0}^{\infty}{\frac{x^n}{n!}} = 1 + \frac{x}{1!} + \frac{x^2}{2!} + \frac{x^3}{3!} + ...\]
\[ \cos{x} = \sum_{n = 0}^{\infty}{\frac{(-1)^n}{(2n)!}}x^{2n} = 1 - \frac{x^2}{2!} + \frac{x^4}{4!} - \frac{x^6}{6!} + ...\]
\[ \sin{x} = \sum_{n = 0}^{\infty}{\frac{(-1)^{n}}{(2n + 1)!}x^{2n + 1}} = \frac{x}{1!} - \frac{x^3}{3!} + \frac{x^5}{5!} - \frac{x^7}{7!} + ... \]
\rl Вызначым у камплекснай плоскасці $e^{z}$, $\sin{z}$ і $\cos{z}$ гэтай сумай. Зараз паспрабуем падставіць у экспаненту $z = i\phi$
\begin{align*}
    e^{i\phi} = \sum_{n = 0}^{\infty}{\frac{(i\phi)^n}{n!}} = \sum_{n = 0}^{\infty}{\frac{(i\phi)^{2n}}{(2n)!}} + \sum_{n = 0}^{\infty}{\frac{(i\phi)^{2n + 1}}{(2n + 1)!}} = \\
    = \sum_{n = 0}^{\infty}{\frac{(-1)^{n}}{(2n)!}\phi^{2n}} + i\sum_{n = 0}^{\infty}{\frac{(-1)^{n}}{(2n + 1)!}\phi^{2n + 1}} = \\
    = \cos{\phi} + i \sin{\phi}
\end{align*}
\rl Такім чынам падстаўляя сінусы і косінусы замест эквівалентных ім радоў атрымліваем формулу Эйлера
\[ e^{i\phi} = \cos{\phi} + i\sin{\phi} \]
\rl З гэтага можна зрабіць вывад, што $e^{i\phi}$ будзе кропкай на акружнасьці даўжынёй 1, радыус да якой будзе ўтвараць вугал $\phi$ супраць гадзіннай стрэлкі
\section{Падрабязьней пра вектара}
\rl Але гэта ж толькі адзіны спосаб задаць вектара. Ёсьць яшчэ другі: мы можам апісаць вектар выкарыстоўвая яго даўжыню і вугал наклону.
Хай гэты вугал $\phi$ будзе супраць гадзіннай стрэлкі. \\
\rl Тады вектар задаецца наступнай парай: $(r, \phi)$, дзе $r$ \--- радыус(даўжыня) вектара.
З формулы Эйлера няцяжка вывесьці выгляд у камплекснай натацыі
\[ e^{i\phi} = \cos{\phi} + i\sin{\phi} \]
\[ (r, \phi) \leftrightarrow re^{i\phi} \]
\rl Нагадаю, што $\cos{\phi}$ і $\sin{\phi}$ гэта даўжыня праэкцыі вектара даўжыні 1 з вуглом $\phi$ на вось $OX$ і $OY$ адпаведна, іншымі словамі яго $x$ і $y$
\section{Пераўтварэньне Фурье}
\rl Як ты мабыць памятаеш, Пераўтварэньне Фурье выглядае наступным чынам
\[ \hat{f}(\omega) = \int_{-\infty}^{\infty}{f(t)e^{i\omega t}dt} \]
\begin{center}
    \scriptsize $f$ \--- сігнал, $\omega$ \--- частата, $t$ \--- час
\end{center}
\rl Але ў дыскрэтным варыяньце яно выглядае так
\[ \hat{f}_k = \sum_{n = 0}^{N - 1}{f_n e^{\frac{2\pi k}{N}n}} \]
\[ \hat{f}_k = \sum_{n = 0}^{N - 1}{f_n \omega^{nk}} \]
\begin{center}
    \scriptsize $k$ \--- індэкс частаты, $N$ \--- даўжыня сігналу, $\omega = e^{\frac{2\pi}{N}}$ \--- першаабразны корань з 1, з формулы Эйлера зразумела,
    што корані з 1 разьмешчаны раўнамерна
\end{center}
\rl Абратнае Пераўтварэньне Фурье выглядае наступным чынам
\[ f(t) = \frac{1}{2\pi} \int_{-\infty}^{\infty}{\hat{f}(\omega)e^{-i\omega t} d\omega} \]
\[ f_n = \frac{1}{N} \sum_{k = 0}^{N - 1}{\hat{f}_k e^{-\frac{2\pi k}{N} n}} \]
\[ f_n = \frac{1}{N} \sum_{k = 0}^{N - 1}{\hat{f}_k \omega^{-nk}}\]
\rl k-e Дыскрэтнае Пераўтварэньне Фурье можна разгледжваць у сэнсе падстаўленьня $z = \omega^k$ у мнагачлен $f(z) = \sum_{n = 0}^{N - 1}{f_n z^n}$, а абратнае \--- інтэрпаляцыяй ісходнага мнагачлена,
аднаўленьне па кропкам $\hat{f}_k$
\section{Як гэта працуе?}
\rl Што ж, мы падабраліся да самой суці. Як гэта працуе? \\
\rl Возьмем тыя кропкі(абазначым $k$-ую з іх $z'_k$), якія мы жадаем намаляваць. Адсарціруем іх па вуглу супраць гадзіннай стрэлкі. Прымянім да іх Абратнае Дыскрэтнае Пераўтварэньне Фурье. 
На выхадзе атрымаем каэфіцыенты да мнагачлена $f(z)$, які пры падстанове $z = \omega^k$ дае нам $z'_k$. Будзем па чаразе падстаўляць у $f(z)$ $z = e^{it}$, дзе $t$ \--- гэта час, каб зрабіць кручэньне вектароў
\section{Рэалізацыя}
\rl Кругі і вектара рысуе відэакарта на заднем фоне, таму на тваём ноутбуке адрысоўка можа быць вельмі доўгая. Для кожнага пікселя фону я вылічваю, ці уваходзіць ён у які-небудзь круг або вектар.
Адначасова на галоўным плане я рысую лініі, якія злучваюць кропкі, якія я палічыў раней
\end{document}